\documentclass{beamer}
\mode<presentation>
\usepackage{amsmath}
\usepackage{amssymb}
%\usepackage{advdate}
\usepackage{graphicx}
\graphicspath{{./figs/}}
\usepackage{adjustbox}
\usepackage{subcaption}
\usepackage{enumitem}
\usepackage{multicol}
\usepackage{mathtools}
\usepackage{listings}
\usepackage{url}
\def\UrlBreaks{\do\/\do-}
\usetheme{Boadilla}
\usecolortheme{lily}
\setbeamertemplate{footline}
{
  \leavevmode%
  \hbox{%
  \begin{beamercolorbox}[wd=\paperwidth,ht=2.25ex,dp=1ex,right]{author in head/foot}%
    \insertframenumber{} / \inserttotalframenumber\hspace*{2ex} 
  \end{beamercolorbox}}%
  \vskip0pt%
}
\setbeamertemplate{navigation symbols}{}

\providecommand{\nCr}[2]{\,^{#1}C_{#2}} % nCr
\providecommand{\nPr}[2]{\,^{#1}P_{#2}} % nPr
\providecommand{\mbf}{\mathbf}
\providecommand{\pr}[1]{\ensuremath{\Pr\left(#1\right)}}
\providecommand{\qfunc}[1]{\ensuremath{Q\left(#1\right)}}
\providecommand{\sbrak}[1]{\ensuremath{{}\left[#1\right]}}
\providecommand{\lsbrak}[1]{\ensuremath{{}\left[#1\right.}}
\providecommand{\rsbrak}[1]{\ensuremath{{}\left.#1\right]}}
\providecommand{\brak}[1]{\ensuremath{\left(#1\right)}}
\providecommand{\lbrak}[1]{\ensuremath{\left(#1\right.}}
\providecommand{\rbrak}[1]{\ensuremath{\left.#1\right)}}
\providecommand{\cbrak}[1]{\ensuremath{\left\{#1\right\}}}
\providecommand{\lcbrak}[1]{\ensuremath{\left\{#1\right.}}
\providecommand{\rcbrak}[1]{\ensuremath{\left.#1\right\}}}
\theoremstyle{remark}
\newtheorem{rem}{Remark}
\newcommand{\sgn}{\mathop{\mathrm{sgn}}}
\providecommand{\abs}[1]{\left\vert#1\right\vert}
\providecommand{\res}[1]{\Res\displaylimits_{#1}} 
\providecommand{\norm}[1]{\lVert#1\rVert}
\providecommand{\mtx}[1]{\mathbf{#1}}
\providecommand{\mean}[1]{E\left[ #1 \right]}
\providecommand{\fourier}{\overset{\mathcal{F}}{ \rightleftharpoons}}
%\providecommand{\hilbert}{\overset{\mathcal{H}}{ \rightleftharpoons}}
\providecommand{\system}[1]{\overset{\mathcal{#1}}{ \longleftrightarrow}}
%\providecommand{\system}{\overset{\mathcal{H}}{ \longleftrightarrow}}
	%\newcommand{\solution}[2]{\textbf{Solution:}{#1}}
%\newcommand{\solution}{\noindent \textbf{Solution: }}
\providecommand{\dec}[2]{\ensuremath{\overset{#1}{\underset{#2}{\gtrless}}}}
\newcommand{\myvec}[1]{\ensuremath{\begin{pmatrix}#1\end{pmatrix}}}
\let\vec\mathbf

\lstset{
%language=C,
frame=single, 
breaklines=true,
columns=fullflexible
}

\numberwithin{equation}{section}
\title{12.319}
\author{AI25BTECH11033 - Spoorthi N}
% \maketitle
% \newpage
% \bigskip
\begin{document}
{\let\newpage\relax\maketitle}
\renewcommand{\thefigure}{\theenumi}
\renewcommand{\thetable}{\theenumi}

\textbf{Question}:

Let $M$ be the real vector space of $2 \times 3$ matrices with real entries.  
Let $T : M \to M$ be defined by
\begin{align}
T\!\left( 
\begin{bmatrix}
x_{1} & x_{2} & x_{3} \\
x_{4} & x_{5} & x_{6}
\end{bmatrix}
\right)
=
\begin{pmatrix}
-x_{6} & x_{4} & x_{1} \\
x_{3} & x_{5} & x_{2}
\end{pmatrix}.
\end{align}

The determinant of $T$ is \_\_\_\_ \hfill (MA 2013)

\textbf{Solution}:

We are given the transformation
\begin{align}
T\!\left(
\begin{pmatrix}
x & y & z \\
u & v & w
\end{pmatrix}
\right)
=
\begin{pmatrix}
-w & u & x \\
z & v & y
\end{pmatrix}.
\end{align}

Represent input as a vector;

\begin{align}
\begin{pmatrix}
x & y & z \\
u & v & w
\end{pmatrix}
\quad \longleftrightarrow \quad
\begin{bmatrix}
x \\ y \\ z \\ u \\ v \\ w
\end{bmatrix}.
\end{align}

Transformation on the vector;

\begin{align}
T\!\left(
\begin{pmatrix}
x \\ y \\ z \\ u \\ v \\ w
\end{pmatrix}
\right)
=
\begin{pmatrix}
-w \\ u \\ x \\ z \\ v \\ y
\end{pmatrix}.
\end{align}

Matrix of transformation;

\begin{align}
A =
\begin{pmatrix}
0 & 0 & 0 & 0 & 0 & -1 \\
0 & 0 & 0 & 1 & 0 & 0 \\
1 & 0 & 0 & 0 & 0 & 0 \\
0 & 0 & 1 & 0 & 0 & 0 \\
0 & 0 & 0 & 0 & 1 & 0 \\
0 & 1 & 0 & 0 & 0 & 0
\end{pmatrix}.
\end{align}

Determinant:

\begin{align}
\det(A) = (-1)\times(-1) = 1.
\end{align}

\begin{align}
{\det(T) = 1}
\end{align}

\end{document}
